\chapter{Implantable Venous Access Port}

\section*{Overview}
An implantable venous access port (Port-a-Cath) is a small reservoir and catheter placed under the skin that provides durable central venous access for repeated therapies such as chemotherapy, fluids, or blood sampling. The exact approach, setting, and preparation depend on the indication, anatomy, and the patient's health.

\section*{Alternative Names}
Port-a-Cath; subcutaneous port; implanted port; central venous port

\section*{Principle}
A catheter is inserted into a central vein, typically the internal jugular or subclavian, and tunneled to a subcutaneous reservoir port placed on the chest wall. The port is accessed through the skin with a non-coring needle for infusion or aspiration.

\section*{History}
Implantable ports were developed in the 1980s following the growth of long-term chemotherapy, building on earlier central venous catheter and subcutaneous reservoir technology.

\section*{Why It Is Done}
Performed for patients needing frequent or long-term intravenous therapy, most commonly chemotherapy, total parenteral nutrition, antibiotics, or blood products, when peripheral access is inadequate or undesirable.

\section*{Is This a Primary or Secondary Test or Procedure}
Primary durable vascular access; secondary to peripheral or external central access for long-term use.

\section*{How It Is Performed}
Performed under local anesthesia with or without sedation, usually with ultrasound and fluoroscopy. The vein is accessed percutaneously, the catheter is positioned with its tip at the cavoatrial junction, and the port is placed in a subcutaneous pocket on the chest and connected to the catheter. The incisions are closed.

\section*{Preparation}
Preoperative assessment of coagulation, site examination, and NPO if sedation is planned. Aseptic technique is essential.

\section*{Contraindications and Precautions}
Active bacteremia or infection at the site, and severe coagulopathy are relative contraindications. Precaution: careful technique avoids pneumothorax and arterial puncture.

\section*{Risks}
Pneumothorax, bleeding, infection, catheter thrombosis, pinch-off syndrome, catheter migration, and wound complications.

\section*{Aftercare and Recovery}
The port is flushed and used for therapy. Patients may shower after a few days; the port remains accessible for years with regular flushing.

\section*{Outcome and Follow-up}
Fluoroscopy confirms correct tip position. Patency and absence of complications are verified before use.

\section*{Expected Results}
A functional port that delivers and aspirates blood easily, with no infection, thrombosis, or mechanical complication.

\section*{Follow-up}
Flushing every 4-6 weeks when unused, and prompt evaluation of any swelling, pain, or fever for infection or thrombosis.

\section*{When to Seek Medical Care}
Follow the procedure team's specific instructions. Seek urgent medical assessment for severe or worsening pain, uncontrolled bleeding, fever or signs of infection, trouble breathing, chest pain, fainting, new weakness or confusion, inability to urinate, or any rapidly worsening symptom. The appropriate warning signs depend on the procedure and the patient's medical condition.

\section*{References}
\begin{enumerate}
  \item MedlinePlus. Port-a-Cath. U.S. National Library of Medicine; 2025.
  \item Current Medical Diagnosis and Treatment. McGraw-Hill; 2025.
\end{enumerate}

\clearpage
